\section{Cotangent and exterior restriction on the original geometric twist}

\begin{definition}[The original native cotangent restriction]
\label{def:cartier-geometric-native-cotangent-restriction}
\leanok
\uses{def:cartier-geometric-twisted-open-restrictions}
\lean{EnvelopeProperty.Cartier.GeometricCotangentRestrictions.nativeCotangentRestriction,EnvelopeProperty.Cartier.GeometricCotangentRestrictions.nativeCotangentRestriction_D,EnvelopeProperty.Cartier.GeometricCotangentRestrictions.rawCotangentIso,EnvelopeProperty.Cartier.GeometricCotangentRestrictions.rawCotangentIso_D}
For an inclusion $U\subset V$, the original native twist map
$\tau:T_V\to T_U$ induces the universal cotangent map
$\Omega_{T_V/k}\to\tau_*\Omega_{T_U/k}$, sending $dt$ to $d\tau(t)$.
For an original affine open $U$, the original section-ring isomorphism
$\sigma_U:\Gamma(Y,U^{(p)})\simeq T_U$ induces the universal
cotangent isomorphism, sending every $da$ to $d\sigma_U(a)$.
\end{definition}

\begin{theorem}[The entire raw cotangent restriction square]
\label{thm:cartier-geometric-raw-cotangent-restriction}
\leanok
\uses{def:cartier-geometric-native-cotangent-restriction,thm:cartier-geometric-original-sections-iso-naturality}
\lean{EnvelopeProperty.Cartier.GeometricCotangentRestrictions.rawCotangentIso_restrict}
For every inclusion $U\subset V$ of original affine opens, the
original raw cotangent restriction commutes with the original
$\sigma$-induced cotangent isomorphisms and the native cotangent
restriction on every cotangent element.
\end{theorem}
\begin{proof}
View both composites as linear maps over $\Gamma(Y,V^{(p)})$,
using the proved whole original section-ring square to identify
their target scalar actions. The module of relative differentials
is generated by the universal differentials $da$. On $da$, the
first composite has value $d\tau(\sigma_V(a))$. Naturality of the
original universal derivation shows that the second has value
$d\sigma_U(a|_{U^{(p)}})$. The whole original section-ring square
identifies these two arguments. Universal uniqueness for relative
differentials gives equality of the two entire linear maps.
\end{proof}

\begin{lemma}[Original source-unit restriction]
\label{lem:cartier-geometric-original-source-unit-restriction}
\leanok
\uses{def:cartier-geometric-twisted-open-restrictions}
\lean{EnvelopeProperty.Cartier.GeometricCotangentRestrictions.cotangent_unit_restrict,EnvelopeProperty.Cartier.GeometricCotangentRestrictions.exterior_unit_restrict,EnvelopeProperty.Cartier.GeometricCotangentRestrictions.twistedUnitIso_restrict}
The original cotangent sheafification unit and, in every degree,
the original exterior sheafification unit commute with the
original inclusion $U^{(p)}\subset V^{(p)}$. Whenever honest etale
tuples give the established affine unit isomorphisms,
those same isomorphisms satisfy this entire source-unit square.
\end{lemma}
\begin{proof}
Each unit is a morphism of the original presheaves of modules.
Its naturality identity evaluated at an arbitrary raw section is
exactly the asserted square. The established affine unit
isomorphism has forward map equal to the original unit on every
raw section, so substitution gives the same identity for it.
\end{proof}

\begin{theorem}[The entire raw exterior restriction square]
\label{thm:cartier-geometric-raw-exterior-restriction}
\leanok
\uses{thm:cartier-geometric-raw-cotangent-restriction,thm:cartier-geometric-original-sections-iso-naturality}
\lean{EnvelopeProperty.Cartier.GeometricCotangentRestrictions.nativeExteriorRestriction,EnvelopeProperty.Cartier.GeometricCotangentRestrictions.nativeExteriorRestriction_wedge,EnvelopeProperty.Cartier.GeometricCotangentRestrictions.rawDomainIso_restrict}
In every degree $q$, the original raw exterior restriction commutes
with the original $\sigma$-induced exterior isomorphisms and the
native exterior restriction. This includes $q=0$.
\end{theorem}
\begin{proof}
Use the whole original section-ring square to regard both
composites as linear maps over the larger geometric section ring.
Pure wedges of $q$ cotangent elements span the exterior power.
Both universal exterior maps send such a wedge to the wedge of
the corresponding cotangent images. The proved entire cotangent
restriction square identifies each entry of those two wedges,
so the composites agree on a spanning family and hence everywhere.
For $q=0$ the same spanning family contains the empty wedge and
the argument applies unchanged.
\end{proof}
